\documentclass[twocolumn,prl,superscriptaddress,nofootinbib]{revtex4-2}

\usepackage{amsmath,amssymb}
\usepackage{amsthm}
\usepackage{booktabs}
\usepackage{graphicx}
\usepackage{hyperref}

\newtheorem{theorem}{Theorem}

\newcommand{\R}{\mathbb{R}}
\newcommand{\C}{\mathbb{C}}
\newcommand{\Z}{\mathbb{Z}}
\newcommand{\Q}{\mathbb{Q}}
\newcommand{\CK}{\mathrm{CK}\text{-}31}
\newcommand{\Norm}{\mathrm{N}}
\newcommand{\Tr}{\mathrm{Tr}}
\newcommand{\braket}[2]{\langle#1|#2\rangle}

\begin{document}

\title{The arithmetic of Kochen--Specker uncolorability\\in dimension three}

\author{Michael Kernaghan}
\affiliation{Pacific Quantum Systems, Vancouver, Canada}

\date{\today}

\begin{abstract}
Kochen--Specker (KS) sets in dimension three are usually studied through their
orthogonality graphs.  We show that their existence is instead governed by the
arithmetic of the ring supplying their coordinates.  Over the alphabet
$\{0,\pm1,\pm x\}$, orthogonality demands a vanishing three-term sum, and every
such \emph{cancellation identity} can be enumerated: each is an algebraic
constraint on~$x$, and each solution is a KS-supporting ring.  Three mechanisms
occur---$|x|^2=2$, a vanishing sum of roots of unity, and the minimal polynomial
of a cubic unit---while generators of squared modulus ${\geq}\,3$ that are not
roots of unity yield triads but no uncolorability.  KS sets therefore occupy
isolated algebraic islands; we exhibit eight, with certified pool minima
$31 < 33 < 43 < 52 < 58$.  For cyclotomic alphabets the criterion is a theorem:
the $n$th-root-of-unity ray family in~$\C^3$ is KS-uncolorable iff $6 \mid n$.
\end{abstract}

\maketitle

\textit{Introduction.---}The Kochen--Specker (KS) theorem~\cite{KochenSpecker1967}
forbids assigning context-independent definite values to quantum observables (for
a review see Ref.~\cite{Budroni2022}).  It is proved by exhibiting a finite set of
rays admitting no $\{0,1\}$ coloring in which orthogonal rays are never both~1 and
every orthogonal triad carries exactly one~1.  Dimension three is the sharpest
case.  The smallest known set is the Conway--Kochen 31-ray configuration
($\CK$), reported by Peres~\cite{Peres1993}; the best lower bound is~24,
established independently by SAT methods with algebraic realizability
checking~\cite{LiBrightGanesh2024,KirchwegarPeitlSzeider2023}, improving the
earlier bound of~22~\cite{UijlenWesterbaan2016}.  The intervening gap of seven
rays is a central open problem, and Trandafir and Cabello~\cite{TrandafirCabello2025}
have since proved $\CK$ rigid and conjectured it optimal.

Searches for smaller sets have been combinatorial---enumerate candidate
orthogonality hypergraphs, then test geometric realizability~\cite{LiBrightGanesh2024,KirchwegarPeitlSzeider2023}---or
have fixed a coordinate alphabet by hand~\cite{Peres1991,Pavicic2004}.  We take
the opposite route.  Two classic constructions, $\CK$ over $\{0,\pm1,\pm2\}$ and
Peres-33 over $\{0,\pm1,\pm\sqrt2\}$, are displayed on the same cube
diagram~\cite{Peres1993} and share the coordinate template $\{0,\pm1,\pm x\}$.
They differ only in~$x$.  This invites a question whose answer is arithmetic:
\emph{which $x$ support a KS set?}

We find that a single local condition on the coordinate ring controls the answer,
that the condition can be enumerated exhaustively for two-symbol alphabets, and
that for cyclotomic alphabets it is a theorem.  KS sets do not populate the space
of algebraic parameters densely; they sit on isolated islands, which we
enumerate.

\textit{The cancellation criterion.---}Let $\mathcal{A} \subset \C$ be a finite
coordinate alphabet and let $S(\mathcal{A}) = \{[v] : v \in
\mathcal{A}^3\setminus\{0\}\}$ be the induced ray set, $[v]$ denoting the
projective class of~$v$.  Rays $v,w$ are orthogonal when
$\sum_k \bar v_k w_k = 0$: a vanishing sum of at most three terms drawn from the
product set $\bar{\mathcal{A}}\cdot\mathcal{A}$.  We call such a relation a
\emph{cancellation identity}.  Without one there are no orthogonal triads, and
without triads there is no KS set, so the search for KS-supporting alphabets is
the search for cancellation identities---and, as it turns out, for identities of
low norm.

All coordinates are represented as exact algebraic numbers and every
orthogonality is certified to vanish exactly, with no floating-point tolerance.
Colorability is decided by SAT (Glucose4~\cite{AudemardSimon2018}): one Boolean
per ray, an exactly-one clause group per triad, an at-most-one clause per
orthogonal pair; UNSAT means KS-uncolorable.

For $\mathcal{A} = \{0,\pm1,\pm x\}$ the product set is
$\{0,\pm1,\pm x,\pm\bar x,\pm|x|^2\}$, and the primitive three-term vanishing
sums admit a complete enumeration (Table~\ref{tab:zerosum}), certified with an
SMT solver~\cite{deMouraBjorner2008}.  Each row is an algebraic constraint
on~$x$; each solution is a coordinate ring.

\begin{table}[b]
\caption{\label{tab:zerosum}All primitive three-term zero-sums available to the
product set of $\mathcal{A}=\{0,\pm1,\pm x\}$.  Every row is a constraint
on~$x$; every solution yielding uncolorability is one of the islands of
Table~\ref{tab:islands}.  The enumeration applies to the raw alphabet, before
cross-product completion.}
\begin{ruledtabular}
\begin{tabular}{lll}
Zero-sum & Constraint on $x$ & Outcome \\
\colrule
$1+1-x$            & $x=2$                    & Integer \\
$1+1-|x|^2$        & $|x|^2=2$                & Peres, $\Z[\sqrt{-2}]$, Heeg.-7 \\
$1+x-|x|^2$        & $|x|^2=1+x$              & Golden \\
$1+x+\bar x$       & $\Tr(x)=-1$, $\Norm(x)=1$ & Eisenstein \\
$1-x-\bar x$       & $\Tr(x)=+1$              & colorable\footnote{Holds for
every imaginary quadratic generator with $d \equiv 1 \pmod 4$.  It produces
all-nonzero triads, but these decouple from the one-zero triads and the ray set
stays colorable; the obstruction is SMT-certified for all integer generator
norms $N \geq 3$.} \\
$2|x|^2-1$         & $|x|^2=\tfrac12$         & $\equiv$ Peres (rescaled) \\
$2x-|x|^2$         & $|x|^2=2x$               & $\equiv$ Integer \\
\end{tabular}
\end{ruledtabular}
\end{table}

Three mechanisms appear among the rows that produce uncolorability.
\emph{Modulus-2 cancellation} ($|x|^2=2$) covers $\sqrt2$, $\sqrt{-2}$, $1+i$ and
$\alpha=(1+\sqrt{-7})/2$; the integer case $1+1=2$ is the degenerate instance in
which $x=2$ and the identity arises additively.  \emph{Phase cancellation} is the
vanishing sum of unit-modulus terms $1+\omega+\omega^2=0$, $\omega=e^{2\pi i/3}$.
In Galois terms the two are clean: modulus-2 is
$|\Norm_{K/\Q}(x)|=2$, and phase is $\Norm_{K/\Q}(x)=1$ with
$\Tr_{K/\Q}(x)=-1$, which characterizes the roots of $\Phi_3$.

A third mechanism exists, and we record it as a correction.  Earlier versions of
this work asserted that modulus-2 and phase cancellation exhaust the possibilities.
That thesis was falsifiable---a KS set over an alphabet whose product set admits
neither identity would refute it---and Gonzalez communicated exactly such a set:
47~rays over $\Q(\psi)$, where $\psi \approx 1.4656$ is the supergolden ratio,
$\psi^3=\psi^2+1$.  We reproduced it in exact arithmetic (47 rays, 107 orthogonal
pairs, 29 triads, ray-critical).  Its raw product set $\{0,\pm1,\pm\psi,\pm\psi^2\}$
contains no element of squared modulus~2; $\Q(\psi)$ is a real field of odd degree
and so contains no root of unity but~$\pm1$; and no additive $1+1-2=0$ cancellation
occurs anywhere in the critical set.  The operative relation is the minimal
polynomial itself, giving a third mechanism: \emph{minimal-polynomial
cancellation}, in which the defining relation of a cubic algebraic unit supplies
the vanishing sum directly.

The negative half of the criterion is equally sharp and, across our survey,
without exception: an alphabet whose generator has $|x|^2 \geq 3$ and is not a
root of unity produces orthogonal pairs and triads, but never enough interlocking
triads for uncolorability.  Note that $|\psi|^2 \approx 2.148$ lies below~3, so
the supergolden island respects this bound while violating the two-mechanism
classification.

One caveat is structural.  Closing a ray set under cross products---for each
orthogonal pair $v \perp w$, adjoin $[v \times w]$, and in~$\C^3$ the Hermitian
cofactor analogue---can introduce coordinates absent from $\mathcal{A}$, and hence
cancellation identities the raw alphabet does not support.  The classification is
therefore properly one of \emph{completed} coordinate algebras.  The golden ratio
is the case in point: $\{0,\pm1,\pm\varphi\}$ is colorable, but completion
generates $1/\varphi$ and the completed pool is not.

\textit{Cyclotomic alphabets.---}For one family the criterion is not a survey
finding but a theorem.

\begin{theorem}
The ray set generated in $\C^3$ by the alphabet $\{0\}\cup\{\zeta_n^k\}$,
$\zeta_n=e^{2\pi i/n}$, is KS-uncolorable if and only if $6 \mid n$.
\end{theorem}

\noindent\emph{Sketch.}  Inner products in $\C^3$ involve at most three terms, so
only short vanishing sums matter.  A three-term sum $\zeta^a+\zeta^b+\zeta^c=0$
exists iff $3 \mid n$: dividing by $\zeta^a$ and requiring $|{-}1-\alpha|=1$ forces
$\cos(2\pi p/n)=-1/2$.  A two-term sum requires $\zeta^{b-a}=-1$, hence $2\mid n$.
Sufficiency then follows by embedding: if $6\mid n$ the sixth roots sit inside the
$n$th roots, carrying the Eisenstein 33-ray set~\cite{Cabello2025simplest} with
them.  Necessity splits into three cases.  With neither $2\mid n$ nor $3\mid n$ the
only orthogonalities are axis-complementary and a single triad survives.  With
$2\mid n$ alone, no pair of all-nonzero rays is orthogonal; every triad has the
form $\{v_1,v_2,e_k\}$, each coordinate plane decomposes into a perfect matching,
and an explicit coloring follows.  With $3\mid n$ alone---the principal
case---each all-nonzero ray $v$ has exactly two projectively distinct orthogonal
partners, $w_\pm$, and since $\braket{w_+}{w_-}=1+\omega+\omega^2=0$ the triple
$\{v,w_+,w_-\}$ is the unique triad containing~$v$.  For odd~$n$ no ray with one
zero coordinate joins any triad.  The triads are therefore isolated, the coloring
problem decomposes into independent satisfiable subproblems, and the pool is
colorable.\hfill$\square$

The condition restates cleanly in Galois terms: $\mathrm{Gal}(\Q(\zeta_n)/\Q)
\cong (\Z/n\Z)^*$ must surject onto $\Z/6\Z$, which holds iff $6\mid n$.  The
prime~6 is not a coincidence of this family.  Cortez, Morales and
Reyes~\cite{CortezMoralesReyes2022} show that $M_3(\Z[1/6])_\mathrm{sym}$ is the
minimal ring extension of~$\Z$ admitting no algebraic hidden states.  Their
invariant $N(S)=\mathrm{lcm}\{\|v\|^2\}$, which fixes the ring $\Z[1/N(S)]$ in
which the projectors $P_v$ live, takes the value~6 exactly on the Eisenstein
island, 12 on Peres, $\Z[\sqrt{-2}]$ and Heegner-7, and 30 on~$\CK$.  Every
island needs the primes 2 and~3; only the integer island additionally needs~5.

\textit{The islands.---}We call a KS-supporting pair (ring, alphabet) an
\emph{algebraic island} because it is isolated in parameter space: perturbing the
generator off its exact algebraic value destroys uncolorability immediately, and
no continuous path joins one island to another.  A sweep of the one-parameter
family $\{0,\pm1,\pm\cos\theta,\pm\sin\theta\}$ over 164 angles in $(0,\pi/2]$
makes this concrete---exactly five angles are uncolorable, all reducing to known
islands, and moving $1^\circ$ from any of them restores colorability.

Table~\ref{tab:islands} lists the eight islands found across real quadratic fields
$\Q(\sqrt d)$ for non-square $d \leq 30$, all nine imaginary quadratic fields of
class number~1, cyclotomic fields for $n \leq 60$, and selected cubic and quartic
extensions.  Each pool minimum is certified by an OCUS (optimal constrained
unsatisfiable subset) procedure that exhausts all smaller candidates within the
pool---a proof within each finite pool, not a heuristic upper bound.

\begin{table*}[t]
\caption{\label{tab:islands}The eight algebraic islands.  ``Pool'' is the ray count
after completion where completion is required; ``Min'' is the OCUS-certified
minimum KS subset within that pool and $|\mathcal{B}|$ its basis count.  Peres,
$\Z[\sqrt{-2}]$ and the Gaussian realization $\{0,\pm1,\pm i,\pm(1{+}i)\}$ share a
single orthogonality graph (VF2-verified), so the eight islands carry seven
distinct graph types.}
\begin{ruledtabular}
\begin{tabular}{llllcccl}
Island & Ring & Alphabet & Identity & Mechanism & Pool & Min & $|\mathcal{B}|$ \\
\colrule
Integer        & $\Z$                     & $\{0,\pm1,\pm2\}$                    & $1+1=2$                 & modulus-2          & 49  & \textbf{31} & 17 \\
Peres          & $\Z[\sqrt2]$             & $\{0,\pm1,\pm\sqrt2\}$               & $(\sqrt2)^2=2$          & modulus-2          & 49  & 33 & 16 \\
Eisenstein     & $\Z[\omega]$             & $\{0,\pm1,\pm\omega,\pm\bar\omega\}$ & $1+\omega+\omega^2=0$   & phase              & 57  & 33 & 14 \\
$\Z[\sqrt{-2}]$& $\Z[\sqrt{-2}]$          & $\{0,\pm1,\pm\sqrt{-2}\}$            & $|\sqrt{-2}|^2=2$       & modulus-2          & 49  & 33 & 16 \\
Heegner-7      & $\Z[\tfrac{1+\sqrt{-7}}2]$ & $\{0,\pm1,\pm\alpha,\pm\bar\alpha\}$ & $\alpha\bar\alpha=2$  & modulus-2          & 145 & 43 & 23 \\
Golden         & $\Z[\varphi]$            & $\{0,\pm1,\pm\varphi\}$              & $\varphi^2=\varphi+1$   & modulus-2 (compl.) & 205 & 52 & 25 \\
Cubic          & $\Z[\sqrt[3]2]$          & $\{0,\pm1,\pm\sqrt[3]2,\pm\sqrt[3]4\}$ & $\sqrt[3]2\cdot\sqrt[3]4=2$ & modulus-2 (compl.) & 361 & 58 & 26 \\
Supergolden    & $\Z[\psi]$               & $\{0,\pm1,\pm\psi\}$                 & $\psi^3=\psi^2+1$       & min.\ polynomial   & 157 & $\leq 47$ & --- \\
\end{tabular}
\end{ruledtabular}
\end{table*}

Two features of the table carry the argument.  First, the hierarchy tracks the
algebraic complexity of the identity rather than the ambient field: the integer
identity produces the densest triad network (26 triads from 49 rays) and the
smallest set, while indirect and minimal-polynomial cancellations produce the
sparsest and the largest.  Second, ambient field is genuinely not the invariant.
The Peres and $\Z[\sqrt{-2}]$ islands have identical ray, pair and triad counts,
and their orthogonality graphs are isomorphic at both the pair and the triad
level; the Gaussian alphabet $\{0,\pm1,\pm i,\pm(1{+}i)\}$ reproduces the same
graph again.  All three realize $|x|^2=2$ in different rings.  Of 28 two-generator
alphabets tested, the 17 reaching minimum~33 fall into exactly two graph classes:
13 containing a modulus-2 generator give the Peres graph, 4 containing $\omega$
give the Eisenstein graph.  The cancellation identity, not the field, fixes the
graph.

The imaginary quadratic case is the cleanest illustration.  Among the nine
Heegner numbers $d=1,2,3,7,11,19,43,67,163$, exactly four support KS sets---and
they are exactly the four whose ring-of-integers generator has
$|\alpha|^2=2$ ($d=1,2,7$) or is a primitive cube root of unity ($d=3$).  The
remaining five have $|\alpha|^2 = 3,5,11,17,41$ and are colorable.  Extending
beyond unique factorization changes nothing: across 26 imaginary quadratic fields
of class number 2, 3 and~4, every field with $|\alpha|^2 \geq 8$ is colorable at
both basic and enriched alphabets.  Class number does not predict uncolorability;
generator norm does.

\textit{New configurations.---}Two islands yield configurations we have not found
in the published catalogues.  The Heegner-7 ring gives a 43-vector set in
23~bases with degree sequence $\{4{:}16,\,5{:}20,\,6{:}4,\,8{:}2,\,10{:}1\}$,
structurally richer than either the Eisenstein 33-set (three degree types) or the
Peres 33-set (two).  The golden ratio field gives a 52-vector set that is
invisible to raw-alphabet search: $\{0,\pm1,\pm\varphi\}$ yields only 10 triads
and is colorable, and uncolorability appears only after completion expands the
pool to 205 rays.  That $\varphi$ works where $\sqrt5$ does not---the same field,
a different generator---is the sharpest single demonstration that the alphabet,
not the field, is the carrier of the structure.

Both new sets are infinitesimally rigid in~$\C^3$, tested by comparing the null
space of the Jacobian of the normalization-and-orthogonality system against the
symmetry dimension $n+8$.  Four of the six non-cubic islands are rigid on this
test.  The two exceptions, Peres and $\Z[\sqrt{-2}]$, each carry exactly one
deformation, and the flex is finite rather than infinitesimal: it integrates to
the three-phase family of Gould and Aravind~\cite{GouldAravind2010}, whose gauge
reduction is a single circle~\cite{FloresGordillo2026} carrying the two islands at
antipodal parameter values.  Since Trandafir and Cabello~\cite{TrandafirCabello2025}
classify known constructions into dimension-lifting and concatenation families and
show that neither yields rigid sets, the algebraic-alphabet method falls outside
both.

\textit{Operational consequences.---}The substrate is not merely a bookkeeping
device.  Every bipartite perfect quantum strategy---the maximum algebraic
violation of a Bell inequality under projective measurement---contains a KS
set~\cite{Cabello2024bipartite}, and face nonsignaling correlations, full Bell
nonlocality, all-versus-nothing proofs and pseudotelepathy are manifestations of
the same structure~\cite{Cabello2024equivalences}.  The islands are therefore the
algebraic substrates of all perfect strategies in dimension three.  Cabello's
``simplest KS set''~\cite{Cabello2025simplest}, built from the Weyl--Heisenberg
group over $\Z[\omega]$, is the smallest member of our Eisenstein island; that an
independently motivated construction lands on an island rather than between two is
evidence the classification tracks real structure.

Computing bipartite KS input counts by SAT across the islands, following the
optimal conversion of Trandafir and Cabello~\cite{TrandafirCabello2025optimal},
shows that the arithmetic, not the ray count, sets the cost of the resulting Bell
scenario.
Eisenstein achieves $5\times9$ (exhaustive), Peres and $\Z[\sqrt{-2}]$ both
$7\times9$, while $\CK$---with the fewest rays of any island---needs at least
8~Alice inputs and reaches only $8\times9$.  Heegner-7 and Golden cost
$\leq 9\times12$ and $\leq 12\times13$.  $\CK$ is thus strictly dominated as a
strategy engine by sets with more rays but fewer contexts.  The trade-off is
genuinely three-way once Cabello--Severini--Winter invariants~\cite{CabelloSeveriniWinter2014}
are included: $\CK$ alone retains a nontrivial sandwich
$\alpha=11<\vartheta=11.71<\alpha^*=12$, while the strongest contextual advantage
of any pool, $\vartheta/\alpha=1.118$, belongs to Heegner-7 and is destroyed by
minimization---the rays that carry it participate in no basis, so cardinality
minimization discards them first.  Three islands sit on the resulting Pareto
frontier and none dominates on every axis.

\textit{Discussion.---}The picture that emerges is that contextuality in dimension
three is not a continuous geometric phenomenon but a discrete arithmetic one.  A
coordinate ring either carries a low-norm cancellation identity or it does not,
and the answer decides not only whether a KS set exists but how small it can be,
how rigid it is, and how cheap a Bell scenario it generates.

Two limits should be stated plainly.  First, the survey is finite.  The
enumeration in Table~\ref{tab:zerosum} is complete for raw two-symbol alphabets,
and the $6\mid n$ result is a theorem, but the eight islands are what a systematic
search over low-degree extensions found, not a classification of all number
fields.  Multi-generator alphabets and extensions of degree five and above remain
the principal gaps.  Second, the completeness of the three-mechanism
classification is open, and the history here is cautionary: the two-mechanism
version survived eight revisions before the supergolden island refuted it, and it
was refuted by an alphabet class we had surveyed but not surveyed far enough.

The constraints are nonetheless directly usable by the lower-bound program.  The
gap between the constructive bound~31 and the realizability bound~24 is being
attacked by SMT search over candidate hypergraphs~\cite{LiBrightGanesh2024,KirchwegarPeitlSzeider2023},
and a candidate graph must ultimately be realized over some coordinate ring.  The
cancellation classification, the discrete jump from the 13-ray pool of
$\{0,\pm1\}$ to the 49-ray pool of $\{0,\pm1,\pm2\}$, and the observation that
every tested construction reaching 31~rays produces the same orthogonality
graph---71 edges, 17 bases, degree sequence $(3^4,4^{14},5^8,6^3,8^2)$---are all
algebraic filters on that search.  If $\CK$ is optimal, the reason is likely to be
arithmetic: $1+1=2$ is the cheapest cancellation identity there is.

\begin{acknowledgments}
The author thanks Ad\'an Cabello for the connection between the Eisenstein island
and the ``simplest KS set''; Manuel Flores Gordillo for identifying a
constraint-ordering error in the rigidity computation; and Jorge Lucas Gonzalez
for the supergolden construction that refuted the two-mechanism thesis, together
with a verifier exact enough to be checked rather than believed.  Computational
work, proof construction and manuscript preparation were carried out with a large
language model (Claude, Anthropic) under the author's direction, with adversarial
cross-review by a second model (GPT, OpenAI); all computational claims are
reproducible from deterministically seeded scripts, and all claims about published
results were checked against the cited papers.  The author directed the research,
validated results against the literature, and accepts full responsibility for the
scientific content.  Code and data:
\url{https://github.com/michaelkernaghan/contextuality}.
\end{acknowledgments}

\begin{thebibliography}{99}

\bibitem{KochenSpecker1967}
S.~Kochen and E.~P.~Specker, J. Math. Mech. \textbf{17}, 59 (1967).

\bibitem{Budroni2022}
C.~Budroni, A.~Cabello, O.~G\"uhne, M.~Kleinmann, and J.-{\AA}.~Larsson,
Rev. Mod. Phys. \textbf{94}, 045007 (2022).

\bibitem{Peres1993}
A.~Peres, \textit{Quantum Theory: Concepts and Methods}
(Kluwer, Dordrecht, 1993).

\bibitem{LiBrightGanesh2024}
Z.~Li, C.~Bright, and V.~Ganesh, in \textit{Proc. 33rd Int. Joint Conf. on
Artificial Intelligence (IJCAI-24)} (IJCAI, 2024), pp.~1898--1906.

\bibitem{KirchwegarPeitlSzeider2023}
M.~Kirchweger, T.~Peitl, and S.~Szeider, in \textit{Proc. AAAI-23} (2023),
pp.~4119--4126.

\bibitem{UijlenWesterbaan2016}
S.~Uijlen and B.~Westerbaan, New Gener. Comput. \textbf{34}, 3 (2016).

\bibitem{TrandafirCabello2025}
S.~Trandafir and A.~Cabello, Phys. Rev. A \textbf{111}, 052204 (2025);
arXiv:2501.11640.

\bibitem{Peres1991}
A.~Peres, J. Phys. A \textbf{24}, L175 (1991).

\bibitem{Pavicic2004}
M.~Pavi\v{c}i\'c, J.-P.~Merlet, B.~McKay, and N.~D.~Megill,
J. Phys. A \textbf{38}, 1577 (2005).

\bibitem{AudemardSimon2018}
G.~Audemard and L.~Simon, Int. J. Artif. Intell. Tools \textbf{27}, 1840001 (2018).

\bibitem{deMouraBjorner2008}
L.~de~Moura and N.~Bj{\o}rner, in \textit{TACAS 2008}, LNCS \textbf{4963}
(Springer, 2008), p.~337.

\bibitem{Cabello2025simplest}
A.~Cabello, arXiv:2508.07335.

\bibitem{CortezMoralesReyes2022}
I.~Cortez, C.~Morales, and M.~Reyes, arXiv:2211.13216.

\bibitem{GouldAravind2010}
E.~Gould and P.~K.~Aravind, Found. Phys. \textbf{40}, 1096 (2010).

\bibitem{FloresGordillo2026}
M.~Flores Gordillo, preprint (2026), Zenodo doi:10.5281/zenodo.21488474.

\bibitem{Cabello2024bipartite}
A.~Cabello, Phys. Rev. Lett. \textbf{134}, 010201 (2025); arXiv:2311.17735.

\bibitem{TrandafirCabello2025optimal}
S.~Trandafir and A.~Cabello, Phys. Rev. A \textbf{111}, 022408 (2025);
arXiv:2410.17470.

\bibitem{Cabello2024equivalences}
Y.~Liu, H.~Y.~Chung, E.~Z.~Cruzeiro, J.~R.~Gonzales-Ureta, R.~Ramanathan, and
A.~Cabello, Phys. Rev. Research \textbf{6}, L042035 (2024); arXiv:2310.10600.

\bibitem{CabelloSeveriniWinter2014}
A.~Cabello, S.~Severini, and A.~Winter, Phys. Rev. Lett. \textbf{112}, 040401 (2014).

\end{thebibliography}

\end{document}
